\documentclass[../m00-session-01.tex]{subfiles}
\begin{document}
\TopicStandaloneTitle{01}{Set the workspace boundary}{Foundations: Python, Cloud and a reproducible workspace}
\TopicDivider{01}{Set the workspace boundary}{Separate the project, credentials, data and documentation from the start.}

\begin{frame}{What this topic establishes}
  \begin{LearningObjectives}{By the end of this topic, learners can}
    \begin{itemize}
      \item describe the role of Python, IAM, S3 and Git in a data-science workflow;
      \item create a minimal local project structure that can be reproduced by a teammate; and
      \item distinguish credentials, code and data artefacts by where they should live.
    \end{itemize}
  \end{LearningObjectives}
\end{frame}

\begin{frame}{The smallest useful data-science workspace}
  \TwoColumnLayout{
    \begin{ConceptBox}{Local project}
      \begin{itemize}
        \item versioned Python environment
        \item \InlineCode{src/} for reusable code
        \item \InlineCode{notebooks/} for exploration
        \item \InlineCode{README.md} for the next learner
      \end{itemize}
    \end{ConceptBox}
  }{
    \begin{ConceptBox}{AWS account boundary}
      \begin{itemize}
        \item IAM role grants least-privilege access
        \item S3 stores durable data and artefacts
        \item CloudShell offers a browser-based terminal
        \item never commit keys or tokens to Git
      \end{itemize}
    \end{ConceptBox}
  }
\end{frame}
\end{document}
